\section{GİRİŞ}

Trafik güvenliği, özellikle yoğun nüfuslu kentsel alanlarda giderek artan bir toplumsal sorun olmaya devam etmektedir. Yaya geçitlerinde araçların geçiş önceliğini ihlal etmesi, ülkemizde ve dünyada ciddi yaralanma ve ölüm vakalarına yol açmaktadır. Türkiye Emniyet Genel Müdürlüğü verilerine göre trafik kazalarının önemli bir bölümü yaya geçidi ihlallerinden kaynaklanmakta; bu durum hem güvenlik hem de hukuki açıdan kapsamlı önlemler alınmasını zorunlu kılmaktadır. Mevcut denetim sistemleri büyük ölçüde manuel yöntemlere ya da pahalı altyapı yatırımları gerektiren sabit kamera sistemlerine dayanmakta, bu nedenle geniş çaplı yaygınlaşma sağlanamamaktadır.

Son yıllarda bilgisayarlı görü ve derin öğrenme alanındaki hızlı gelişmeler, gömülü sistemlerin hesaplama kapasitesindeki artışla birleşerek "Edge-AI" adı verilen yeni bir yaklaşımın önünü açmıştır. Edge-AI, görüntü işleme ve çıkarım işlemlerinin merkezi bir sunucuya aktarılmadan doğrudan uç cihazda ya da yakın birimlerde gerçekleştirilmesine olanak tanımaktadır. Bu yaklaşım; gecikme süresini azaltmakta, bant genişliği tüketimini düşürmekte ve sistemin daha düşük maliyetle konuşlandırılmasını mümkün kılmaktadır. YOLOv8 gibi hafifletilmiş gerçek zamanlı nesne tespit modelleri, standart bir dizüstü bilgisayarda dahi saniyede onlarca kare işleyebilmekte; bu durum Edge-AI tabanlı trafik izleme sistemlerini pratik ve uygulanabilir hale getirmektedir.

Bu proje kapsamında, düşük maliyetli bir ESP32-S3 kamera modülü ile Python tabanlı işleme motorunun birleşiminden oluşan SmartWalk sistemi tasarlanmış ve geliştirilmiştir. Sistem; WiFi üzerinden alınan MJPEG video akışını gerçek zamanlı olarak işlemekte, Canny+Hough algoritması ile yaya geçidi bölgesini otomatik tespit etmekte ve YOLOv8n modeli aracılığıyla araç ve yaya tespiti yapmaktadır. İhlal anında anlık fotoğraf kaydedilmekte, Telegram üzerinden bildirim gönderilmekte ve Firebase Realtime Database'e log kaydı düşülmektedir. Projenin temel katkısı, pahalı özel donanım gerektirmeden erişilebilir bileşenlerle kurulabilecek, genişletilebilir ve açık kaynaklı bir ihlal tespit mimarisi ortaya koymaktır.

